\documentclass[12pt,a4paper]{article}

\usepackage[T1]{fontenc}
\usepackage[utf8]{inputenc}
\usepackage[brazil]{babel}
\usepackage{graphicx}
\usepackage{geometry}
\usepackage{listings}
\usepackage{listingsutf8}
\usepackage{xcolor}
\usepackage{setspace}
\usepackage{float}
\usepackage{tikz}
\usepackage{tabularx}

\usetikzlibrary{positioning, shapes.geometric}
\onehalfspacing
\geometry{margin=2.5cm}

\lstset{
    language=SQL,
    basicstyle=\ttfamily\small,
    keywordstyle=\color{blue},
    commentstyle=\color{gray},
    stringstyle=\color{red},
    breaklines=true,
    showstringspaces=false
}

\newcommand{\montarinicio}[3]{
    % =========================
    % CAPA
    % =========================
    \begin{titlepage}
        \centering

        {\large Faculdade Anhanguera}\\[0.5cm]
        {\large Tecnólogo em DevOps}\\[3cm]

        {\Large \textbf{#1}}\\[2cm]

        {\large Gabriel Blank Stift Mousquer}\\[1cm]

        {\large Ivoti/RS}\\
        {\large 2026}

    \end{titlepage}

    % =========================
    % FOLHA DE ROSTO
    % =========================
    \begin{titlepage}
        \centering

        {\large Faculdade Anhanguera}\\[0.5cm]
        {\large Tecnólogo em DevOps}\\[2cm]

        {\Large \textbf{#1}}\\[1.5cm]

        \begin{flushright}
            Trabalho apresentado à disciplina de \\
            \textbf{#2} \\
            Professor: #3 \\
        \end{flushright}

        \vfill

        {\large Gabriel Blank Stift Mousquer}\\
        Matrícula: 2024201126\\

        \vfill

        {\large Ivoti/RS}\\
        {\large 2026}

    \end{titlepage}

    % =========================
    % SEÇÃO: TOC
    % =========================
    \tableofcontents
    \newpage
}


\begin{document}

\montarinicio
    {ELEMENTOS DO MODELO ENTIDADE-RELACIONEMENTO}
    {Modelagem de Dados}
    {Mariana Karina Miglionari De Souza}

% =========================
% SEÇÃO: IDENTIFICAÇÃO DE ELEMENTOS EM DER
% =========================
\section{Sistema de Biblioteca: Identificação de Elementos em DER}

\subsection{Entidades identificadas}
\begin{itemize}
\item Livro
\item Leitor
\item Empréstimo
\end{itemize}

\subsection{Atributos da entidade LIVRO}
\begin{itemize}
\item ISBN
\item Título
\item Autor
\item Ano de publicação
\end{itemize}

\subsection{Chave primária de LIVRO e justificativa}

A chave primária é o \textbf{ISBN}, pois é um atributo único para cada livro no domínio, garantindo a identificação unívoca de cada instância da entidade Livro.

\subsection{Atributo composto em LEITOR}

O atributo composto é \textbf{Endereço}, que pode ser decomposto em:
\begin{itemize}
\item Rua
\item Número
\item Bairro
\item Cidade
\item CEP
\end{itemize}

\subsection{Relacionamentos e cardinalidades}

\begin{itemize}
\item LEITOR -- EMPRÉSTIMO: \textbf{1 : N}
\begin{itemize}
\item Um leitor pode realizar vários empréstimos
\item Cada empréstimo pertence a um único leitor
\end{itemize}

\item LIVRO -- EMPRÉSTIMO: \textbf{1 : N}
\begin{itemize}
\item Um livro pode estar associado a vários empréstimos
\item Cada empréstimo refere-se a um único livro
\end{itemize}
\end{itemize}

\subsection{Chaves estrangeiras na entidade EMPRÉSTIMO}

\begin{itemize}
\item Matrícula do Leitor (referência à entidade Leitor)
\item ISBN do Livro (referência à entidade Livro)
\end{itemize}

% =========================
% SEÇÃO: DESENHO DO DER
% =========================
\section{Sistema de Clínica Médica: Desenho do DER}

\begin{figure}[h]
\centering
\begin{tikzpicture}[node distance=3cm]

% ENTIDADES (retângulos)
\node (consulta) [draw, diamond] {Consulta};
\node (paciente) [draw, rectangle, left of=consulta, xshift=-2cm] {Paciente};
\node (medico) [draw, rectangle, right of=consulta, xshift=2cm] {Médico};

% ATRIBUTOS PACIENTE
\node (cpf) [draw, ellipse, above of=paciente] {\underline{CPF}};
\node (nomeP) [draw, ellipse, above left of=paciente] {Nome};
\node (telP) [draw, ellipse, double, above right of=paciente] {Telefone};
\node (dataNasc) [draw, ellipse, below of=paciente] {\shortstack{Data de\\nascimento}};

% ATRIBUTOS MÉDICO
\node (crm) [draw, ellipse, above of=medico] {\underline{CRM}};
\node (nomeM) [draw, ellipse, above right of=medico] {Nome};
\node (telM) [draw, ellipse, above left of=medico] {Telefone};
\node (esp) [draw, ellipse, below of=medico] {Especialidade};

% ATRIBUTOS CONSULTA
\node (data) [draw, rectangle, below left of=consulta] {Data};
\node (valor) [draw, rectangle, below of=consulta] {Valor};
\node (horario) [draw, rectangle, below right of=consulta] {Horário};

% LIGAÇÕES PACIENTE
\draw (paciente) -- (cpf);
\draw (paciente) -- (nomeP);
\draw (paciente) -- (dataNasc);
\draw (paciente) -- (telP);

% LIGAÇÕES MÉDICO
\draw (medico) -- (crm);
\draw (medico) -- (nomeM);
\draw (medico) -- (esp);
\draw (medico) -- (telM);

% LIGAÇÕES CONSULTA
\draw[dashed] (consulta) -- (data);
\draw[dashed] (consulta) -- (valor);
\draw[dashed] (consulta) -- (horario);

% RELACIONAMENTOS
\draw (paciente) -- node[midway, fill=white]{(1,N)} (consulta);
\draw (medico) -- node[midway, fill=white]{(1,N)} (consulta);

\end{tikzpicture}
\end{figure}

\subsection{Entidades identificadas}

Foram identificadas \textbf{3 entidades} no cenário:

\begin{itemize}
\item Paciente
\item Médico
\item Consulta
\end{itemize}

\subsection{Cardinalidade do relacionamento MÉDICO--CONSULTA}

O relacionamento entre MÉDICO e CONSULTA é \textbf{1 : N}, pois:

\begin{itemize}
\item Um médico pode realizar várias consultas
\item Cada consulta é realizada por um único médico
\end{itemize}

\subsection{Atributo multivalorado em PACIENTE (telefone)}

O atributo \textbf{telefone} é multivalorado porque:

\begin{itemize}
\item Um paciente pode possuir mais de um número de telefone
\item Não há limitação de unicidade para esse atributo no domínio
\item Portanto, o atributo não é atômico, podendo assumir múltiplos valores para uma mesma instância de paciente
\end{itemize}

% =========================
% SEÇÃO: INTEGRIDADE REFERENCIAL
% =========================
\section{Sistema de Vendas: Integridade referencial}

\subsection{Violação de integridade referencial na tabela PEDIDO}

Sim, há violação de integridade referencial. O pedido de número \textbf{1003} contém o valor \textbf{CPF\_cliente = 444.444.444-44}, porém este CPF não existe na tabela \textbf{CLIENTE}. Portanto, trata-se de uma chave estrangeira sem referência válida.

\subsection{Uso de CPF\_cliente nulo no pedido 1004}

A existência de \textbf{CPF\_cliente = NULL} no pedido 1004 só é permitida se a chave estrangeira tiver sido definida sem restrição de obrigatoriedade (NOT NULL). No modelo relacional, o valor NULL em uma chave estrangeira significa que o relacionamento entre as tabelas não está estabelecido para aquele registro específico. Isso implica que o pedido existe independentemente da associação imediata com um cliente válido. Entretanto, do ponto de vista de modelagem de domínio, isso é incomum, pois um pedido normalmente é um objeto dependente de um cliente. Assim, em um modelo bem estruturado, o atributo CPF\_cliente geralmente deveria ser obrigatório, tornando o uso de NULL semanticamente questionável.

\subsection{Exclusão do cliente com CPF 111.111.111-11}

O CPF do cliente está sendo referenciado como chave estrangeira na tabela PEDIDO (pedido 1001). Assim, a exclusão violaria a integridade referencial. Portanto, a alternativa correta é:

\textbf{Não pode excluir pois há pedidos deste cliente}

\subsection{Regra de integridade referencial para exclusão automática de pedidos}

Para garante que, ao excluir um cliente, todos os pedidos associados a ele sejam automaticamente removidos, mantendo a consistência do banco de dados, a regra de integridade referencial que devemos usar é:

\textbf{ON DELETE CASCADE}

\end{document}
